\chapterrule
\section{前馈网络与输出头：两种不同的特征变换}
\label{intro:10b-ffn-output-head}
\chaptercompass{读取上下文后，怎样变换特征并预测词表？}{FFN 在每行做共享非线性变换，输出头映射为 logits。}{$X:T\times d_{\rm model}\to F:T\times d_{\rm model}\to L:T\times|\mathcal V|$}{$W_{\rm up},W_{\rm down},W_{\rm vocab}$（及相应偏置）}
\modelpipelinefocus{4}{FFN}

\begin{idea}
\noindent\emph{这一章的核心问题}\par

我们刚刚看到完整 Transformer Block 中包含一个 FFN。把早先学过的两层 \conceptlink{mlp}{MLP} 放进整个序列后，输入、输出尺寸怎么选择？
又为什么模型最后还要再乘一次矩阵才能输出词表概率？
同样是线性层，处在计算链上的位置不同，承担的任务就不同。
\end{idea}

\subsection{FFN 是“每个 Token 各做一次”的非线性计算}
\phantomsection\label{concept:ffn}

设有 $T$ 个 Token，隐藏状态为
$X\in\R^{T\times d_{\rm model}}$。
给每个 Token 应用\textbf{相同参数}的两层前馈网络（FFN）：
\begin{equation}
F_{\rm mid}=\relu(XW_{\rm up}+\vect b_{\rm up}),\qquad
F_{\rm out}=F_{\rm mid}W_{\rm down}+\vect b_{\rm down}.
\label{eq:ffn-batched}
\end{equation}
这里的偏置向量按行广播，意味着每一行加上同一组数。
两次矩阵乘法本身只有
\[
XW_{\rm up}=Z,\qquad F_{\rm mid}W_{\rm down}=F_{\rm out}.
\]
形状放在旁边的小表，更容易逐行核对：
\begin{center}\small
\tagpdfsetup{table/header-rows={1}}
\begin{tabularx}{\linewidth}{@{}>{\raggedright\arraybackslash}p{25mm}>{\raggedright\arraybackslash}X@{}}
\toprule
步骤 & 左矩阵 $\times$ 右矩阵 $\to$ 结果\\\midrule
上投影 & $(T\times d_{\rm model})\,(d_{\rm model}\times d_{\rm ff})\to T\times d_{\rm ff}$\\
下投影 & $(T\times d_{\rm ff})\,(d_{\rm ff}\times d_{\rm model})\to T\times d_{\rm model}$\\
\bottomrule
\end{tabularx}
\end{center}
$d_{\rm ff}$ 是 FFN 的中间宽度。它可以大于 $d_{\rm model}$，
但\textbf{不要求固定为四倍}。

\subsection{手算一个两 Token 的 FFN}

为简化例子，设偏置为零，$T=2$、
$d_{\rm model}=2$、$d_{\rm ff}=3$：
\[
X=\begin{bmatrix}1&2\\0&1\end{bmatrix},\quad
W_{\rm up}=\begin{bmatrix}1&0&1\\0&1&-1\end{bmatrix},
\quad
W_{\rm down}=\begin{bmatrix}2&0\\1&1\\1&1\end{bmatrix}.
\]
先上投影：
\[
XW_{\rm up}
=\begin{bmatrix}1&2&-1\\0&1&-1\end{bmatrix}.
\]
逐元素 ReLU：
\[
F_{\rm mid}=\begin{bmatrix}1&2&0\\0&1&0\end{bmatrix}.
\]
再下投影：
\begin{equation}
F_{\rm out}=F_{\rm mid}W_{\rm down}=
\boxed{\begin{bmatrix}4&2\\1&1\end{bmatrix}}.
\end{equation}
这次第一条 Token 从 $(1,2)$ 变成 $(4,2)$，
第二条从 $(0,1)$ 变成 $(1,1)$。两行都发生了变换，
两行的变换独立完成，但都使用同一组 FFN 参数。

\textbf{观察关键事实：}第一行 FFN 不需要第二行参与就可以计算；
第二行也一样。\textbf{标准 FFN 负责逐 Token 处理特征，
\conceptlink{attention}{Attention} 才负责在 Token 之间混合信息。}
如果 FFN 的输入来自 Attention，它处理的向量就已经融合了上下文信息。

\subsection{输出头不是 FFN 的“下投影”}
\phantomsection\label{concept:output-head}

FFN 的下投影将隐藏宽度还原为 $d_{\rm model}$，
是为了继续堆叠后续 Block；
而\textbf{词表输出头}要把隐藏状态映射到候选 Token 分数：
\begin{equation}
Z_{\rm vocab}=F_{\rm out}W_{\rm vocab}+\vect b_{\rm vocab},
\quad
W_{\rm vocab}\in\R^{d_{\rm model}\times|\mathcal V|}.
\end{equation}
假设 $d_{\rm model}=2$，词表有五个 Token，
则 $W_{\rm vocab}$ 是 $2\times5$，
每个 Token 产生五个 \conceptlink{logits}{logits}，即还未经过 Softmax 的原始候选分数。
这些数字是可正可负的 logits，经 Softmax 变成概率。

\begin{center}
\begin{tikzpicture}[>=Latex,node distance=5mm,
 n/.style={draw=linegray,fill=paper,rounded corners=3pt,minimum width=105mm,minimum height=10mm,align=center,font=\small}]
\node[n] (x) {$T\times d_{\rm model}$：每个 Token 的隐藏表示};
\node[n,below=of x] (up) {FFN 上投影 $T\times d_{\rm ff}$，激活，再下投影};
\node[n,below=of up] (y) {$T\times d_{\rm model}$：供下一个 Block 使用};
\node[n,below=of y] (out) {词表输出头 $T\times|\mathcal V|$：预测候选分数};
\draw[->,blue,thick] (x)--(up);
\draw[->,blue,thick] (up)--(y);
\draw[->,teal,thick] (y)--(out);
\end{tikzpicture}
\end{center}

\subsection{再与 Attention 对照一次}

\begin{center}
\small
\renewcommand{\arraystretch}{1.3}
\tagpdfsetup{table/header-rows={1}}
\begin{tabularx}{\linewidth}{@{}p{24mm}p{31mm}X@{}}
\toprule
组件 & 通常处理的轴 & 主要功能\\
\midrule
Attention & Token 与 Token & 动态决定读取哪些位置的信息\\
FFN & 每个 Token 的特征 & 非线性变换当前表示\\
词表输出头 & 特征到词表类别 & 计算下一 Token 的候选分数\\
\bottomrule
\end{tabularx}
\end{center}



\chapterchecktasks{10b-ffn-output-head}{正文基础三题}{
\item \textbf{定位：}给出 FFN 与词表输出头各自的输入／输出单位：它们是隐藏特征变换还是词表类别打分？
\item \textbf{手算：}若 $T=4,d=6,d_{\rm ff}=18$，求 FFN 上投影和下投影的矩阵形状，以及整个 FFN 的输出形状。
\item \textbf{找错：}有人说“FFN 把每行的特征变成词表概率，因此无需再接输出头或 Softmax”。指出两个漏掉的步骤。
}
\chapterrecallanswer{FFN 在各 Token 内逐行改变隐藏特征；输出头才把隐藏宽度映射为词表 logits，Softmax 再变概率。}



\begin{conclusion}
同一隐藏宽度进、同一宽度出，是堆叠 Block 的便利条件；
词表投影则是把内部表示变成预测任务答案的接口。
两者都包含矩阵乘法，但在网络中承担不同的任务。
下一章不再孤立地算两个 Token；我们将沿用前文的\textbf{三维、三个位置}统一参数，把位置、两个 Norm、Attention、残差、FFN 与输出头全部接起来。
\end{conclusion}
